\documentclass[10pt]{article}
\usepackage[T1]{fontenc}
\usepackage[utf8]{inputenc}
\usepackage[french]{babel}
\usepackage{amsmath,amssymb,mathrsfs,mathpazo}
\usepackage{geometry}
\geometry{a4paper,margin=2cm}
\usepackage{graphicx,array,tabularx,tabularray,colortbl,multirow,multicol,booktabs,enumitem,pifont,ulem,fancybox,diagbox,cancel}
\usepackage{pstricks,pstricks-add,pst-plot,pst-tree,pst-node,pst-text,pst-func,pst-math,pst-bspline,pst-3dplot,pst-eucl,pst-solides3d}
\let\pscalculate\undefined
\usepackage{pst-bezier}
\usepackage{tikz,pgfplots,tkz-tab}
\usetikzlibrary{arrows,arrows.meta,patterns,calc,positioning,trees,decorations.markings,shapes.geometric,shapes.misc,shapes.arrows,intersections,decorations.pathmorphing,decorations.pathreplacing,backgrounds,shadows,babel}
\pgfplotsset{compat=1.18}
\usepgfplotslibrary{fillbetween}
\usepackage[np]{numprint}
\usepackage{listings,makecell,pas-tableur,twoopt}
\let\barre\undefined
\usepackage[pstricks]{tablvar}
\definecolor{darkgreen}{RGB}{0,100,0}
\providecommand{\red}{\color{red}}
\providecommand{\blue}{\color{blue}}
\providecommand{\green}{\color{green}}
\providecommand{\black}{\color{black}}
\providecommand{\white}{\color{white}}
\providecommand{\gray}{\color{gray}}
\providecommand{\cyan}{\color{cyan}}
\providecommand{\magenta}{\color{magenta}}
\providecommand{\yellow}{\color{yellow}}
\providecommand{\violet}{\color{violet}}
\providecommand{\orange}{\color{orange}}
\providecommand{\euro}{\texteuro}
\providecommand{\EUR}{\texteuro}
\providecommand{\ieme}{\textsuperscript{e}}
\providecommand{\degres}{\ensuremath{^\circ}}
\providecommand{\degree}{\ensuremath{^\circ}}
\providecommand{\vv}[1]{\overrightarrow{#1}}
\providecommand{\Coord}[1]{#1}
\providecommand{\decosix}{}
\providecommand{\starredbullet}{\ensuremath{\bullet}}
\pagestyle{empty}
\setlength{\parindent}{0pt}
\usepackage{adjustbox}\DeclareUnicodeCharacter{2236}{:}\DeclareUnicodeCharacter{2190}{\ensuremath{\leftarrow}}
\providecommand{\R}{\mathbb{R}}
\providecommand{\N}{\mathbb{N}}
\providecommand{\Z}{\mathbb{Z}}
\providecommand{\Q}{\mathbb{Q}}
\providecommand{\D}{\mathbb{D}}
\providecommand{\C}{\mathbb{C}}
\providecommand{\e}{{\mathrm{e}}}
\providecommand{\vect}[1]{\overrightarrow{#1}}
\providecommand{\vectt}[1]{\overrightarrow{\mathrm{#1}}}
\providecommand{\vvt}[1]{\overrightarrow{\mathrm{#1}}}
\providecommand{\np}[2][]{{#2~#1}}
\providecommand{\vv}[1]{{\overrightarrow{#1}}}
\renewcommand{\d}{{\mathrm{d}}}
\providecommand{\degres}{{^\circ}}
\providecommand{\degre}{{^\circ}}
\providecommand{\blue}{\color{blue}}
\providecommand{\red}{\color{red}}
\providecommand{\enskip}{\quad}
\providecommand{\llbracket}{\lbrack\!\lbrack}
\providecommand{\rrbracket}{\rbrack\!\rbrack}
\providecommand{\euro}{{\text{€}}}
\providecommand{\no}{{\text{n°}}}
\providecommand{\interval}[2]{{[#1;#2]}}
\providecommand{\intervFO}[2]{{[#1;#2[}}
\providecommand{\dotfill}{\ldots}

\providecommand{\ds}{\displaystyle}
\providecommand{\textding}[1]{#1}
\providecommand{\Oij}{\(\mathrm{O};\vec{\imath},\vec{\jmath}\)}
\providecommand{\Oijk}{\(\mathrm{O};\vec{\imath},\vec{\jmath},\vec{k}\)}
\providecommand{\Ouv}{\(\mathrm{O};\vec{u},\vec{v}\)}
\providecommand{\intervFF}[2]{\ensuremath{[#1;#2]}}
\renewcommand{\euro}{\ifmmode\text{\texteuro}\else\texteuro\fi}
\DeclareUnicodeCharacter{00B0}{\ifmmode^{\circ}\else\textdegree\fi}
\DeclareUnicodeCharacter{20AC}{\ifmmode\text{\texteuro}\else\texteuro\fi}
% Énoncés : Éducation nationale. Transcription et corrigés : APMEP, Denis Vergès et contributeurs. Découpage et classement : CoopMaths, Sébastien Lozano, Éric Elter et contributeurs. Adaptation de la présentation : Maths Explorer.
% CC BY-SA : https://coopmaths.fr/www/about/
% https://coopmaths.fr/alea/static/bac/2022/tex/bac_2022_05_sujet1_madagascar_1.tex
\begin{document}
\noindent\textit{APMEP / CoopMaths. CC BY-SA. Présentation adaptée par Maths Explorer.}\par\medskip
\medskip

\begin{tabularx}{\linewidth}{|X|}\hline
\textbf{Principaux domaines abordés :} Probabilités\\ \hline
\end{tabularx}

\medskip

Dans une station de ski, il existe deux types de forfait selon l'âge du skieur :

\begin{itemize}
\item un forfait JUNIOR pour les personnes de moins de 25 ans ;
\item un forfait SENIOR pour les autres.
\end{itemize}

Par ailleurs, un usager peut choisir, en plus du forfait correspondant à son âge l'\emph{option  coupe-file} qui permet d'écourter le temps d'attente aux remontées mécaniques.

On admet que:
\begin{itemize}
\item[$\bullet~~$] 20\,\% des skieurs ont un forfait JUNIOR ;
\item[$\bullet~~$] 80\,\% des skieurs ont un forfait SENIOR ;
\item[$\bullet~~$] parmi les skieurs ayant un forfait JUNIOR, 6\,\% choisissent l'option coupe-file ;
\item[$\bullet~~$] parmi les skieurs ayant un forfait SENIOR, 12,5\,\% choisissent l'option coupe-file.
\end{itemize}

On interroge un skieur au hasard et on considère les évènements :

\begin{itemize}
\item[$\bullet~~$] $J$ : \og le skieur a un forfait JUNIOR \fg ;
\item[$\bullet~~$] $C$ : \og le skieur choisit l'option coupe-file \fg.
\end{itemize}

\medskip

\emph{Les deux parties peuvent être traitées de manière indépendante}

\bigskip

\textbf{Partie A}

\medskip

\begin{enumerate}
\item Traduire la situation par un arbre pondéré.
\item Calculer la probabilité $P(J \cap C)$.
\item Démontrer que la probabilité que le skieur choisisse l'option coupe-file est égale à $0,112$.
\item Le skieur a choisi l'option coupe-file. Quelle est la probabilité qu'il s'agisse d'un
skieur ayant un forfait SENIOR ? Arrondir le résultat à $10^{-3}$.
\item Est-il vrai que les personnes de moins de vingt-cinq ans représentent moins de 15\,\% des skieurs ayant choisi l'option coupe-file ? Expliquer.
\end{enumerate}

\bigskip

\textbf{Partie B}

\medskip

On rappelle que la probabilité qu'un skieur choisisse l'option coupe-file est égale à 0,112.

On considère un échantillon de $30$ skieurs choisis au hasard.

Soit $X$ la variable aléatoire qui compte le nombre des skieurs de l'échantillon ayant
choisi t'option coupe-file.

\medskip

\begin{enumerate}
\item On admet que la variable aléatoire $X$ suit une loi binomiale. 

Donner les paramètres de cette loi.
\item Calculer la probabilité qu'au moins un des $30$ skieurs ait choisi l'option coupe-file. Arrondir le résultat à $10^{-3}$.
\item Calculer la probabilité qu'au plus un des $30$ skieurs ait choisi l'option coupe-file. Arrondir le résultat à $10^{-3}$.
\item Calculer l'espérance mathématique de la variable aléatoire $X$.
\end{enumerate}

\bigskip

\end{document}
